\documentclass[titlepage]{article}
\usepackage[margin=0.5in]{geometry}
\usepackage{graphicx}
\usepackage{subcaption}
\usepackage{caption}
\usepackage[normalem]{ulem}
\usepackage{amsfonts}
\usepackage{amsmath,amssymb}
\usepackage{pifont}
\usepackage{amsthm}
\usepackage{listings}
\usepackage{algorithm}
\usepackage{algpseudocode}
\usepackage{float}
\usepackage[table,dvipsnames]{xcolor}
\usepackage{longtable}
\usepackage{booktabs}
\renewcommand{\algorithmicrequire}{\textbf{Input:}}
\renewcommand{\algorithmicensure}{\textbf{Output:}}
\usepackage[authoryear]{natbib}
\usepackage[colorlinks=true, citecolor=blue, urlcolor=blue, linkcolor=blue]{hyperref}
\usepackage[font=small,labelfont=bf]{caption}
\usepackage[subrefformat=parens]{subcaption}
\usepackage[framemethod=TikZ]{mdframed}
\usepackage{dirtree}
\usepackage{tcolorbox}
\usepackage{makecell}

% Define a custom quote box with only a left vertical bar
\newmdenv[
  topline=false,
  rightline=false,
  bottomline=false,
  linewidth=2pt,
  linecolor=gray,
  innerleftmargin=10pt,
  innerrightmargin=10pt,
  innertopmargin=5pt,
  innerbottommargin=5pt,
  backgroundcolor=none
]{sidequote}


\title{Simple Simulation II: Psychometrics in AI Deliberation}
\author{Yashwit Polapragada}

\setlength{\parindent}{0pt}
\captionsetup{justification=centering}

\begin{document}

\maketitle

\tableofcontents

\newpage

\section{Experiment Overview}
\label{sec:taskdescription}
\rule{\textwidth}{0.4pt}

\subsection{Aim}
The aim of the experiment is simulating \textbf{dyadic} civic deliberation on a range of topics, and thereafter judging the deliberative ability of LLMs relative both to their assigned citizen personas, and to the topic itself.\\

Broadly, this experiment explores the evaluation of deliberative quality using \textbf{psychometric theory}, by appending to the deliberation a confidential policy vote for each deliberative agent. While this is done primarily to challenge popular alternative metrics like semantic convergence and DQI\footnote{Susceptible to the quirks of the underlying RLHF chatbot-training of LLMs; one example being the conversational technique such as mirroring.}, there is empirical basis to suggest that deliberation and voting should be combined to optimize democratic outcomes.\\

Specifically, \texttt{Yang and Bachmann 2025} argue that deliberation and voting are complementary rather than conflicting democratic mechanisms. Post-deliberative voting serves as an essential ``fair closure mechanism'' that guarantees equal standing, encourages authentic preference revelation, and preserves legitimate dissent. Furthermore, \texttt{Yang and Bachmann 2025} suggest that administering this ballot privately as a secret vote is essential to prevent social conformity pressure and prevailing-opinion bias.\\

This experiment defaults to 4 currently relevant issues in Singapore, spanning social, economic, and environmental axes\footnote{As will be emphasised later however, the user can configure as many topics as they wish to simulate deliberations for.}.

\subsubsection{Topic 1: School Closures}
\begin{sidequote}
    \emph{Due to declining birth rates, the Ministry of Education (MOE) regularly merges or closes primary and secondary schools. This is a highly emotional topic, pitting practical resource allocation and diverse social mixing against heritage, alumni loyalty, and localized community disruption.}
\end{sidequote}

Reference: \href{https://www.channelnewsasia.com/singapore/nine-pairs-schools-merge-2022-amid-falling-birth-rates-changing-demographics-5650006}{CNA, 2021}.

\subsubsection{Topic 2: Heritage Shophouses vs New MRT Lines}
\begin{sidequote}
    LTA initially announced that three shophouse units (numbers 235, 237, and 239) of the 1924 conserved Ellison Building would be completely demolished to clear the way for underground excavation works, with plans to reconstruct them after the tunnel's completion. This drew strong objections from organizations like the Singapore Heritage Society and ICOMOS Singapore. They argued that tearing down and rebuilding a gazetted structure undermines the core principles of heritage conservation.
\end{sidequote}

Reference:
\href{https://www.straitstimes.com/singapore/tunnel-works-to-affect-heritage-building}{ST, 2016.}

\subsubsection{Topic 3: Foreign PMET Wage Thresholds}
\begin{sidequote}
    This item captures the Ministry of Manpower’s (MOM) long-running policy calibration between maintaining Singapore's status as a regional corporate/MNC hub and addressing citizen PMET anxieties over wage suppression, displacement, and leadership pipelines.
\end{sidequote}

Reference:
\href{https://www.edb.gov.sg/en/news-and-insights/salary-threshold-for-new-employment-pass-applicants-to-be-raised-to-5600-from-2025}{EDB, 2024.}

\subsubsection{Topic 4: Haze Response}
\begin{sidequote}
    Members of Parliament (including WP MP Kenneth Tiong and others) filed parliamentary questions regarding the zero-conviction track record under the THPA since 2014, demanding to know how Singapore can impose real financial accountability on corporate beneficiaries of slash-and-burn clearing without diplomatic paralysis.
\end{sidequote}

Reference:
\href{https://www.straitstimes.com/singapore/why-cant-singapores-anti-haze-law-get-indonesian-companies-to-stop-the-fires}{ST, 2026.}

\subsection{Setup}
\subsubsection{LLM Personas}

The experiment defaults to $N = 4$ different citizen personas:
\begin{itemize}
    \item $\boldsymbol{P}_1$: The Technocratic Pragmatist; Calm, analytical, focused on feasibility, data, and institutional mechanisms.
    \item $\boldsymbol{P}_2$: The Ideological Advocate; Articulate, morally committed to preservation and equity; skeptical of technocratic coldness.
    \item $\boldsymbol{P}_3$: The Disaffected Populist; Emotive, candid, feels sidelined by elites; defensive and prone to scapegoating heuristics.
    \item $\boldsymbol{P}_4$: The Open-Minded Everyman; Plainspoken, humble, constructive; seeks common-sense harmony.
\end{itemize}

The choice of $N = 4$ personas is based on the matrix below crossing \textit{Knowledge} $\times$ \textit{Polarization}. \textbf{In general however, the experiment allows for any variable number of $N$ personas.}

\begin{table}[h!]
\centering
\renewcommand{\arraystretch}{1.4}
\begin{tabular}{r|c|c|}
\multicolumn{1}{c}{} & \multicolumn{1}{c}{\textbf{Low Polarization}} & \multicolumn{1}{c}{\textbf{High Polarization}} \\
\cline{2-3}
\textbf{High Knowledge} & 
\makecell{\textbf{$P_1$: Technocrat}\\Analytical / Feasibility} & 
\makecell{\textbf{$P_2$: Ideologue}\\Moral Purism / Principled} \\
\cline{2-3}
\textbf{Low Knowledge} & 
\makecell{\textbf{$P_4$: Everyman}\\Consensus / Pragmatic} & 
\makecell{\textbf{$P_3$: Populist}\\Anti-Elite / Emotive} \\
\cline{2-3}
\end{tabular}
\caption{2$\times$2 Persona Matrix.}
\label{tab:compact_personas}
\end{table}

\subsection{Dyadic Configurations}

The deliberations are structured around pairs of the above personas. The number and choices of dyadic configurations are user-configurable (with a theoretical maximum of $\frac{N(N-1)}{2}$). The benchmark defaults to $M = 4$ canonical dyad configurations:

\begin{itemize}
    \item $\boldsymbol{D}_1$ (Technocrat vs Ideologue): focused on friction between personal convictions of what governance looks like (\textit{Ideological Friction}).
    \item $\boldsymbol{D}_2$ (Technocrat vs Populist): focused on agent vulnerability due to information asymmetry (\textit{Asymmetric Vulnerability}).
    \item $\boldsymbol{D}_3$ (Ideologue vs Populist): focused on intractable deadlock between two polarized stances (\textit{Polarized Stalemate}).
    \item $\boldsymbol{D}_4$ (Ideologue vs Everyman): focused on the dynamic between articulate advocacy and common-sense grounding (\textit{Rhetorical Dominance}).
\end{itemize}

\subsection{Procedure}
\begin{enumerate}
    \item User prepares a \texttt{data/items/} folder with individual \texttt{itemXX\_\{topic\}\_description.json} files. The user can specify any number of topics, so long as each config is populated with the required information about the topic to be deliberated.
    \item User prepares a \texttt{data/config.json} file and loads-in an API key and/or Base URL in \texttt{.env}.
    \item User executes \texttt{python3 -m src.main --config "path\_to\_config.json"}.
    \item Script programmatically extracts the config details and initializes client for compute service.
    \item For each topic item, a deliberation is conducted between the user-configured dyadic pairs for the user-configured number of turns.
    \item After the agents have concluded all turns, they vote - privately - on a 4-option MCQ for their policy of choice.
    \item Steps 5 and 6 are repeated for the user-configured number of trials, to produce a more representative outcome for the topic and deliberation configuration.
    \item Raw transcripts for each topic are saved in separate JSON files in the folder \texttt{data/expXXX/transcripts/}, with the naming format \texttt{itemXX\_\{topic\}\_transcript.json}\footnote{The experiment-run indexed \texttt{expXXX} folder is dynamically generated.}. The secret ballots are saved in \texttt{data/expXXX/ballots.json}, alongside snapshots of the configuration (\texttt{config.json}) and battery item definitions (\texttt{items/}) used for the run. 
    \item Raw JSON metrics over all topics over all trials are saved in a single file \texttt{data/expXXX/psychometric\_evals.json}.
    \item A separate \texttt{final\_report.md} file summarizing the run outcome is also programmatically generated. 
\end{enumerate}

\subsection{Outcome}

In essence, the contract of the repo is:
\begin{sidequote}
    \textbf{IF} the user specifies the \texttt{data/items/} and \texttt{data/config.json} properly and loads in a valid API key and/or Base URL in \texttt{.env},\\

    \textbf{THEN} the pipeline will run an $N$-turn deliberation between $M$ user-specified dyads of agent personas over $W$ topics, replicated for $T$-trials; each trial ends with a final policy vote, and then the chat histories and voting outcomes of the trials are recorded in topic-specific \texttt{JSON} files. The pipeline will also output a \texttt{psychometric\_evals.json} with a series of psychometric evaluation figures and overall conclusions. All outputs will be saved in an automatically created \& indexed experiment run sub-directory of \texttt{data}.
\end{sidequote}

In general, if the user requests $W$ topics of deliberation between $M$-configurations of dyads to be replicated for $T$-trials, then there will be a total of $W \times M \times T$ deliberations generated, and $W \times M \times T \times 2$ ballots conducted (since each of the 2 agents in each deliberation vote independently).\\

A sample \texttt{itemXX\_\{topic\}\_transcript.json} looks like:

\begin{verbatim}
[
    {
        "item_id": "ITEM_01",
        "item_title": "School Closures and Mergers in Singapore",
        "dyad_id": "D1",
        "agents": [
          "technocrat",
          "ideologue"
        ],
        "trial_index": 0,
        "turns": [
          {
            "turn_index": 0,
            "speaker": "Technocratic Pragmatist (P1)",
            "persona_id": "technocrat",
            "statement": "..."
          },
          {
            "turn_index": 1,
            "speaker": "Ideological Advocate (P2)",
            "persona_id": "ideologue",
            "statement": "..."
          }
        ]
    },
    {
        ....
    },
]
\end{verbatim}

A sample \texttt{ballots.json} looks like:
\begin{verbatim}
[
  {
    "item_id": "ITEM_01",
    "dyad_id": "D1",
    "trial_index": 0,
    "votes": {
      "technocrat": {
        "item_id": "ITEM_01",
        "dyad_id": "D1",
        "trial_index": 0,
        "persona_id": "technocrat",
        "raw_response": "A",
        "displayed_letter": "A",
        "selected_key": "C",
        "is_correct": false,
        "is_valid": true,
        "display_mapping": {
          "A": "C",
          "B": "D",
          "C": "A",
          "D": "B"
        }
      },
      "ideologue": {
        ...
      }
    },
    "consensus_reached": false
  },
  ...
]
\end{verbatim}

A sample \texttt{psychometric\_evals.json} looks like:
\begin{verbatim}
{
  "ctt": {
    "ITEM_01": {
      "item_id": "ITEM_01",
      "difficulty_p": 0.05,
      "discrimination_index": 0.1667,
      "distractor_counts": {
        "A": 11,
        "B": 14,
        "C": 5,
        "D": 10
      },
      "distractor_ratios": {
        "B": 1.2727,
        "C": 0.4545,
        "D": 0.9091
      }
    },
    ...
  },
  "irt": {
    "items": {
      "ITEM_01": {
        "item_id": "ITEM_01",
        "difficulty_b": 1.319,
        "discrimination_a": 2.2792
      },
      ...
    },
    "deliberators": {
      "dyad_theta": {
        "D1": -0.362,
        "D2": -0.2944,
        "D3": -0.201,
        "D4": -0.0422
      },
      "individual_theta": {
        "everyman": 0.1496,
        "ideologue": -0.303,
        "populist": 0.2496,
        "technocrat": -0.1232
      },
      "persona_by_dyad_theta": {
        "everyman": {
          "D4": 0.1496
        },
        ...
      }
    }
  },
  "individual_success_rates": {
    "technocrat": 0.275,
    "ideologue": 0.1833,
    "populist": 0.4,
    "everyman": 0.35
  },
  "insights": {
    "hardest_item": "ITEM_01",
    "easiest_item": "ITEM_04",
    "most_discriminating_item": "ITEM_04",
    "lowest_competence_dyad": "D1",
    "highest_competence_dyad": "D4",
    "average_dcr_by_distractor": {
      "B": 0.9148,
      "C": 0.6856,
      "D": 0.9602
    },
    "most_tempting_distractor": "D"
  }
}
\end{verbatim}


Optionally, the user can enlist an agent to output a \texttt{itemXX\_\{topic\}\_transcript.pdf} for an easier textual analysis or review of the agents' generated arguments.

\subsection{Evaluation}
\subsubsection{Variables}
\textbf{Independent variables}: 
\begin{itemize}
    \item The topic of deliberation.
    \item The citizen persona.
\end{itemize}

\textbf{Dependent variables}: 
The dependent variables are \textbf{psychometric} variables, influenced by Classical Test Theory (CTT) and Item Response Theory (IRT). In this context, the \texttt{item} is the topic of deliberation and the following metrics are evaluated mostly based on the secret ballots that the agents place after deliberation.\\

\underline{CTT Metrics:}
\begin{enumerate}
    \item \textbf{Consensus Rate}: The proportion of the $W \times M$ deliberations (averaged over $T$ trials) where both participants in the dyad reach consensus by independently selecting the right integrative compromise on their secret ballots as would be expected after deliberation.
    \item \textbf{CTT Item Discrimination Index}: Measures an item's ability to differentiate between high-competence and low-competence dyads by taking the difference of the Item Difficulty/Consensus Rate of the top and bottom 27\% of the cohort.
    \item \textbf{Distractor-to-Correct Ratio}: Quantifies the relative cognitive pull of a bad-faith option against a good integrative compromise option across all $W \times M \times 2$ individual secret ballots cast by agents, averaged over $T$ trials.
\end{enumerate}

\underline{IRT Metrics:}
\begin{enumerate}
    \item \textbf{Item Difficulty}: The latent deliberative ability threshold required for a deliberator to have a 50\% probability of adopting the integrative compromise option in the secret ballot.
    \item \textbf{IRT Item Discrimination}: The slope of the Item Characteristic Curve (ICC) at the Item Difficulty, reflecting how sharply the dilemma differentiates between differing ability levels.
    \item \textbf{Latent Deliberative Ability}: The standardized competence trait of a deliberator, estimated under regularized maximum likelihood. 
    \begin{itemize}
        \item Dyadic ability: Joint consensus convergence where both agents in the dyad simultaneously vote for the integrative compromise.
        \item Individual ability: Persona's marginal baseline propensity to vote for the integrative compromise across all topics.
        \item Partner Interaction Effects: Persona's latent ability conditioned on the specific dyadic partner, isolating communicative synergy and polarization shifts.
    \end{itemize}
\end{enumerate}

\subsubsection{Interpretation}
\begin{enumerate}
    \item Consensus Rate: Measured between [0.0, 1.0]; the larger this metric, the stronger the political consensus.
    \item CTT Item Discrimination: Measured between [-1.0, +1.0]; the larger this metric, the greater the ability of the topic to discriminate deliberators.
    \item Distractor-to-Correct Ratio (DCR): Measured between [0.0, $+\infty$]; the larger this metric, the stronger the power of the distractor in out-competing the integrative compromise.
    \item Item Difficulty: Measured between [-4.0, +4.0] under Gaussian prior; the larger this metric, the greater the difficulty of the topic.
    \item IRT Item Discrimination: Measured between [0.1, 4.0]; the larger this metric, the greater the ability of the topic to discriminate deliberators by ability.
    \item Latent Ability: Measured between [-4.0, +4.0]; the larger this metric, the greater the deliberative ability of the agent or dyad.
\end{enumerate}

\subsubsection{Theoretical/Empirical Basis}
The application of psychometric theory to LLM-simulated deliberation was motivated by \texttt{Ke et al. 2026}, which applied CTT/IRT to evaluate LLM responses to assessments meant for students, to judge the quality of such assessments.\\

\texttt{Bächtiger et al. 2022} also suggest that the very popular quantitative metric, DQI, which has been used to judge LLM-deliberative quality in \texttt{Zhang et al. 2026}, may also be reformed using IRT; making deliberative quality a latent instead of absolute variable.

\newpage

\section{Pilot Experiment Specifics}
\label{sec:taskdescription}
\rule{\textwidth}{0.4pt}

\subsection{Repo Architecture}
\dirtree{%
.1 simple-simulation-psychometric-eval-of-deliberation/.
.2 data/.
.3 config.json \DTcomment{Master simulation configuration (models, dyads, items, IRT priors)}.
.3 items/.
.4 item01\_school\_mergers\_description.json \DTcomment{School Mergers policy item}.
.4 item02\_heritage\_mrt\_description.json \DTcomment{Heritage Shophouses vs. MRT Line item}.
.4 item03\_foreign\_talent\_description.json \DTcomment{Foreign PMET Thresholds item}.
.4 item04\_transboundary\_haze\_description.json \DTcomment{Transboundary Haze Enforcement item}.
.3 exp000/ \DTcomment{Indexed directory for experiment artifacts (auto-increments)}.
.4 config.json \DTcomment{Snapshot of master configuration used for this run}.
.4 items/ \DTcomment{Snapshot of standardized battery items evaluated in run}.
.4 transcripts/ \DTcomment{Item-by-item multi-turn deliberation transcripts}.
.4 ballots.json \DTcomment{Post-deliberation secret votes and parsed selections}.
.4 psychometric\_evals.json \DTcomment{Computed CTT ($p, D, \text{DCR}$) and 2PL IRT ($\theta, b, a$) metrics}.
.4 final\_report.md \DTcomment{Synthesized markdown report of benchmark results with Mermaid diagrams}.
.2 docs/.
.3 coding\_guidelines.md \DTcomment{Architectural rules, typing standards, and testing protocols}.
.3 experiment\_details.tex \DTcomment{Comprehensive formal experimental documentation in \LaTeX}.
.3 experiment\_design\_psychometric\_eval\_deliberation.md \DTcomment{Psychometric benchmark evaluation framework}.
.3 prompt\_templates.md \DTcomment{Citizen persona system prompts and ballot templates}.
.2 src/.
.3 main.py \DTcomment{Master orchestrator (\texttt{execute\_simulation\_run}) and CLI entry point}.
.3 deliberation.py \DTcomment{Multi-turn dyadic deliberation chat loop engine}.
.3 ballot.py \DTcomment{Secret ballot administrator with option randomization and robust parser}.
.3 psychometrics.py \DTcomment{Classical Test Theory and 2PL IRT regularized L-BFGS-B estimation engine}.
.3 personas.py \DTcomment{Citizen persona system prompts, traits, and metadata resolution}.
.3 report.py \DTcomment{Synthesis markdown and Mermaid workflow diagram generator}.
.3 schemas.py \DTcomment{TypedDict structural definitions enforcing static type safety}.
.3 llm.py \DTcomment{OpenAI-compatible client with environment credential validation}.
.2 tests/.
.3 test\_ballot\_parser.py \DTcomment{Unit tests for ballot option extraction and presentation randomization}.
.3 test\_psychometrics.py \DTcomment{Unit tests for CTT formulas, 2PL IRT monotonicity, and prior bounds}.
.3 test\_deliberation.py \DTcomment{Unit tests for dialogue turn structure and history assembly}.
.3 test\_pipeline.py \DTcomment{Full end-to-end dry-run pipeline and CLI execution tests}.
.3 test\_model\_availability.py \DTcomment{Integration test verifying configured LLM availability on endpoint}.
.2 pyproject.toml \DTcomment{Pytest, mypy, and build configurations}.
.2 requirements.txt \DTcomment{Pinned runtime dependencies (numpy, scipy, openai, pytest)}.
.2 README.md \DTcomment{Repository documentation, psychometric guide, and execution tutorial}.
}

\subsection{Overall User Configurability}

\underline{\texttt{itemXX\_\{topic\}\_description.json}:}
\begin{itemize}
    \item \texttt{item\_id}
    \item \texttt{title}
    \item \texttt{topic\_prompt} (user prompt)
    \item \texttt{ballot}: \texttt{options}, \texttt{correct\_key}
\end{itemize}

\underline{\texttt{config.json}:}
\begin{itemize}
    \item \texttt{experiment\_metadata}: name, description, version
    \item \texttt{model\_config}: model assignment (mapping the agent model to the desired LLM)\footnote{Example, \texttt{agent\_model}: \texttt{openai/gpt-oss-120b}"}) and model settings (\texttt{temperature, max\_tokens, deliberation\_max\_tokens, ballot\_max\_tokens, seed}).
    \item \texttt{system\_prompts}: agent persona prompts
    \item \texttt{user\_prompts}: deliberation turn signal, secret ballot invitation.
    \item \texttt{deliberation\_config}: turns per session (\texttt{turns\_per\_item}), trials per topic (\texttt{trials\_per\_item}), MCQ option order randomization seed (\texttt{mcq\_option\_order\_seed}), and dyad configurations (\texttt{dyads}).
\end{itemize}

\subsection{Model Config}
\texttt{gpt-oss-120b} was chosen for all persona agents for the trivial reason that this experiment used \href{https://www.together.ai/}{together.ai} for easy access to compute, and this service has a limited selection of larger models that could be used without a dedicated persistent instance. However, there is extensive precedent for the use of general \texttt{gpt}-series models in LLM-social simulation literature.\\

The default \texttt{temperature} of the deliberation agents was set at 0.7 following \textbf{Argyle et al's (2023)} and \textbf{Wang et al's (2025)} methodologies, to promote the diversity of the model's responses and better approximate a population-representative range of outcomes. This does introduce nondeterminism in textual output which complicates experiment reruns in terms of fairly comparing evaluation metrics, which hinge on the specific flow of deliberation. However, as argued by \texttt{Veri 2026}, replicability in AI-deliberation should not be judged by textual identity but by the stability of analytically relevant structural and cognitive metrics under a fully documented configuration. \\

The default \texttt{seed} of the agents was set to \texttt{None}\footnote{Which means the model elects its own random seed at every new point of inference.}, which further boosts outcome variability and mitigates replication covariance, and seed hacking.\\

The \texttt{deliberation\_max\_tokens} (or general \texttt{max\_tokens}) of 512 (3-5 sentences) is appropriate for the level of verbosity we expect from agents in a short deliberation. Likewise, \texttt{ballot\_max\_tokens} is explicitly allocated at 512 tokens to guarantee that hybrid and reasoning models (e.g., Qwen-3.5, DeepSeek-R1) do not exhaust their token limits on internal reasoning traces before emitting their single-letter voting decision. As for input prompts, the 128K context window of \texttt{gpt-oss-120b} more than suffices.

\subsection{Prompts}
\subsubsection{Topic Setting}
The topics were chosen with the same philosophy in mind as the dyad configurations: to promote the representation of diverse tensions that are frequently encountered in local policy discussions.\\

\underline{Topic 1:}
\begin{sidequote}
    Due to falling birth rates, MOE is merging schools with falling enrollment. One participant prioritizes community heritage and alumni identity, while the other prioritizes educational resource efficiency and modern facilities. Deliberate to find common ground.
\end{sidequote}

\underline{Topic 2:}
\begin{sidequote}
    A planned underground MRT transit line intersects a historic district of conserved pre-war shophouses. One participant argues that preserving architectural identity and cultural heritage is non-negotiable, while the other argues that thousands of daily commuters need fast, affordable, and direct public transit connectivity. Deliberate to find common ground.
\end{sidequote}

\underline{Topic 3:}
\begin{sidequote}
    Singapore balances its role as a regional corporate and tech headquarters hub requiring specialized global talent against domestic mid-career PMET anxieties over fair hiring, wage suppression, and leadership pipelines. One participant champions economic openness and global talent access, while the other prioritizes citizen worker protection and knowledge transfer. Deliberate to find common ground.
\end{sidequote}

\underline{Topic 4:}
\begin{sidequote}
    Following the severe 2026 transboundary haze crisis that drove air quality across Singapore into the Very Unhealthy range (PSI $>$ 200) and prompted NEA to issue Section 10 notices under the Transboundary Haze Pollution Act (THPA), intense debate has erupted over the government's response. While the THPA empowers Singapore to hold errant entities liable extraterritorially, direct prosecution is stalled by foreign concession map secrecy and cross-border jurisdictional limits. One participant demands aggressive unilateral legal prosecution and punitive financial retaliation against regional plantation conglomerates, while the other emphasizes diplomatic caution, sovereign non-interference, and cooperative firefighting assistance. Deliberate to find common ground.
\end{sidequote}

\subsubsection{Ballot Setting}

\underline{General Prompt:}
\begin{sidequote}
    You have completed a structured policy deliberation on: \texttt{title}.\\
    Dilemma Context: \texttt{topic\_prompt}.\\ 
    Deliberation Summary/Transcript: \texttt{transcript\_summary}\\
    Options (in randomized order): \texttt{formatted\_options}\\
    Instructions:\\ 
    Select the single resolution that best resolves the policy tension.\\ 
    Respond ONLY with the single letter corresponding to your choice (e.g., A, B, C, or D) and nothing else.
\end{sidequote}

Below are the policy options presented for each topic. \textbf{Before randomization}\footnote{Which means the following order is not necessarily the order of the options presented to the agent in any trial.}, the options are in the order of:
\begin{enumerate}
    \item Integrative Compromise (ideal post-deliberation choice).
    \item Status Quo Inertia; rigid refusal to adapt (distractor 1).
    \item Technocratic Zero-Sum; ruthless efficiency maximization (distractor 2).
    \item External Scapegoating; diverting domestic accountability onto foreign or bureaucratic scapegoats (distractor 3).
\end{enumerate}

The Integrative Compromises for each topic were constructed based on real-world policies and resolutions after the underlying issue had undergone several iterations of public \& parliamentary discussion/feedback. The real-world references are given as footnotes.\\

\underline{Topic 1 Options:}
\begin{enumerate}
    \item Proceed with mergers to optimize facilities and teacher allocation, while preserving alumni heritage through dedicated on-campus archives, co-named buildings, and free shuttle transit for displaced students.\footnote{\href{https://www.channelnewsasia.com/singapore/moe-school-merger-stamford-telok-kurau-pioneer-200791}{CNA, 2021}.}
    \item Immediately freeze all school mergers indefinitely, funding empty classrooms with deficit spending regardless of cohort sizes.
    \item Aggressively shutter all under-enrolled schools within 12 months with no heritage concessions, redirecting all saved funds purely to top-performing institutions.
    \item Halt local mergers and instead mandate foreign student intake quotas to artificially fill classroom vacancies.
\end{enumerate}

\underline{Topic 2 Options:}
\begin{enumerate}
    \item Adopt deep-tunnel micro-bored engineering with foundation underpinning at a higher initial construction cost, combined with adaptive re-use ground entrances embedded within conserved shophouse facades.\footnote{\href{https://www.straitstimes.com/singapore/bigger-portion-of-conserved-ellison-building-to-be-kept}{ST, 2018}.}
    \item Cancel or indefinitely postpone the MRT station in the area to leave all shophouses undisturbed, leaving residents and commuters reliant on congested road transit.
    \item Demolish the historical shophouses along the alignment via compulsory acquisition to execute the cheapest surface-level cut-and-cover construction, treating heritage as secondary to throughput.
    \item Halt the rail extension and blame foreign conservation advocates and tourists for stalling urban progress, deflecting the infrastructure challenge onto external interests.
\end{enumerate}

\underline{Topic 3 Options:}
\begin{enumerate}
    \item Implement a transparent points-based framework (COMPASS) that evaluates candidate salary benchmarks against local PMET percentiles while scoring firms on nationality diversity and local PMET employment share, balancing global talent access with safeguards for a strong local core\footnote{\href{https://www.channelnewsasia.com/singapore/employment-pass-jobs-foreign-workers-mom-budget-2022-debate-parliament-2538966}{CNA, 2022}.}.
    \item Immediately freeze all Employment Pass issuances and cap existing foreign talent renewal quotas, prioritizing status-quo job protection even if multinational corporations relocate regional operations.
    \item Completely abolish salary benchmarks and hiring ratios to let unfettered market supply determine labor costs, dismissing local worker anxieties as unavoidable frictions of globalization.
    \item Mandate invasive public naming-and-shaming of foreign HR managers and blame multinational hiring syndicates entirely for domestic wage stagnation without improving local training pipelines.
\end{enumerate}

\underline{Topic 4 Options:}
\begin{enumerate}
    \item Combine bilateral firefighting assistance and ASEAN satellite data-sharing with domestic regulatory leverage: enforce strict MAS/SGX sustainable finance taxonomies and mandatory supply-chain traceability that cut off Singapore-based bank lending, capital market access, and public procurement for agribusinesses linked to burning\footnote{\href{https://www.straitstimes.com/singapore/why-cant-singapores-anti-haze-law-get-indonesian-companies-to-stop-the-fires}{ST, 2026}.}.
    \item Rely exclusively on diplomatic appeals and routine offers of regional firefighting assistance through ASEAN channels, accepting periodic transboundary haze as an unavoidable meteorological condition without pursuing corporate accountability.
    \item Unilaterally issue extraterritorial arrest warrants and freeze Singapore-based bank accounts of regional plantation executives regardless of diplomatic fallout, prioritizing punitive optics over bilateral environmental cooperation.
    \item Impose blanket retaliatory tariffs on all agricultural imports from neighboring countries and deflect domestic public health obligations by blaming external governments entirely.
\end{enumerate}

\subsubsection{Personas}

\underline{Technocratic Pragmatist ($\boldsymbol{P}_1$):}
\begin{sidequote}
    You are participating in a citizen policy deliberation panel in Singapore.\\
    Role: Technocratic Pragmatist (High Policy Knowledge, Low Affective Polarization)\\[4pt]
    Core characteristics:
    \begin{itemize}
        \item You approach public policy dilemmas with analytical rigor, fiscal awareness, and institutional realism.
        \item You reference policy mechanisms, economic incentives, operational feasibility, and data-driven constraints.
        \item Your tone is calm, professional, and objective. You avoid emotional escalation.
        \item You prioritize structural viability, but you are willing to compromise on tactical details if core policy efficiency and fiscal sustainability are preserved.
    \end{itemize}
\end{sidequote}

\underline{Ideological Advocate ($\boldsymbol{P}_2$):}
\begin{sidequote}
    You are participating in a citizen policy deliberation panel in Singapore.\\
    Role: Ideological Advocate (High Policy Knowledge, High Affective Polarization)\\[4pt]
    Core characteristics:
    \begin{itemize}
        \item You have deep policy knowledge but strong, principled moral commitments (e.g., social equity, cultural heritage, climate justice).
        \item You are deeply skeptical of bureaucratic expediency and technocratic cost-benefit calculations that sacrifice community values.
        \item Your tone is articulate, passionate, and firm.
        \item You vigorously challenge technocratic complacency, though you can recognize solutions that genuinely uphold your core principles.
    \end{itemize}
\end{sidequote}

\underline{Disaffected Populist ($\boldsymbol{P}_3$):}
\begin{sidequote}
    You are participating in a citizen policy deliberation panel in Singapore.\\
    Role: Disaffected Populist (Low Policy Knowledge, High Affective Polarization)\\[4pt]
    Core characteristics:
    \begin{itemize}
        \item You feel everyday Singaporeans are left behind by out-of-touch elites, sterile policies, and rising costs.
        \item You lack interest in technical jargon; you care about bread-and-butter struggles, gut fairness, and perceived double standards.
        \item Your tone is emotive, candid, sometimes skeptical or defensive.
        \item You are prone to status-quo inertia, zero-sum distrust, or scapegoating heuristics unless someone takes your personal anxieties seriously.
    \end{itemize}
\end{sidequote}

\underline{Open-Minded Everyman ($\boldsymbol{P}_4$):}
\begin{sidequote}
    You are participating in a citizen policy deliberation panel in Singapore.\\
    Role: Open-Minded Everyman (Low Policy Knowledge, Low Affective Polarization)\\[4pt]
    Core characteristics:
    \begin{itemize}
        \item You represent the average citizen navigating daily life without deep policy expertise.
        \item You speak plainly, listen actively, and respect differing viewpoints.
        \item Your tone is humble, constructive, and warm.
        \item You look for common-sense, harmonious solutions that help people get along and solve problems practically.
    \end{itemize}
\end{sidequote}

These personas were engineered as inspired by \textbf{Zhuang et al. 2026}'s suggestion that ``role-specific prompts [which] simulate distinct institutional positions and value commitments...[ensure] that deliberations were grounded in plausible conflicts of interest and institutional trade-offs rather than abstract argumentation".

\subsubsection{Turn Signal}
\begin{sidequote}
    You are in a structured policy deliberation with another citizen regarding: \texttt{title}.\\
    Dilemma Context:\texttt{topic\_prompt}.\\
    Conversation History:\texttt{conversation\_history}.\\
    Now it is your turn to speak. In 3-5 sentences, reply directly to the other participant. Express your perspective, respond constructively to their arguments, and work toward a viable policy compromise that balances competing legitimate values.
\end{sidequote}

\subsection{Deliberation Config}

\subsubsection{Number of Turns}
The choice of a fixed $N=4$ total alternating dialogue turns per session (2 turns per agent) is arbitrary for this conservative experiment (both in terms of compute and time).

\subsubsection{Trials per Topic}
The choice of $T=5$ replication trials per deliberation is arbitrary for this conservative experiment.

\subsubsection{MCQ Option Order Seed}
The choice of seed 42 for the randomized ordering of the ballot options presented to the deliberative agent is arbitrary.

\subsubsection{Dyad Configurations}
$M=4$ dyad configurations were chosen based on observing a limited but diverse set of tensions that often crop up in democratic deliberation conducted by heterogeneous groups.\\

\underline{$\boldsymbol{D}_1$: Technocrat vs Ideologue:}\\
This configuration was chosen because it represents impersonal \textit{Ideological Friction} regarding what governance should look like. The Technocrat optimizes for fiscal viability and mechanism feasibility, while the Ideologue optimizes for moral purity, social justice, or heritage continuity.\\

\underline{$\boldsymbol{D}_2$: Technocrat vs Populist:}\\
This configuration was chosen because it represents a \textit{Vulnerability due to Information-Asymmetry}. The Populist may feel vulnerable to being patronized, out-argued, or minimized by the Technocrat's abstract or institutional manner of speaking, possibly stoking defensive backlash.\\

\underline{$\boldsymbol{D}_3$: Ideologue vs Populist:}\\
This configuration was chosen because it represents a \textit{Polarized Stalemate}. When two high-polarization agents interact, neither will concede ground. Without an analytical bridge, the dialogue may reach an ideological impasse where both are likely to retreat into status-quo preservation or scapegoating.\\

\underline{$\boldsymbol{D}_4$: Ideologue vs Everyman:}\\
This configuration was chosen because it represents \textit{Rhetorical Dominance}. The articulate Ideologue is likely to drive the agenda and set the terms of debate, with the consequence of the Everyman either getting steamrolled or managing to anchor the discussion in practical common sense \& underserved concerns.

\newpage{}

\section{Pilot/Demo Experiment Results}
\label{sec:taskdescription}
\rule{\textwidth}{0.4pt}

\newcommand{\cmark}{\textcolor{green!60!black}{\ding{51}}}
\newcommand{\xmark}{\textcolor{red}{\ding{55}}}

\subsection{Evaluation Metrics}
\begin{tcolorbox}[colback=gray!5!white,colframe=gray!75!black,title=Analysis]
\begin{verbatim}
Run A (Control):
  Convergence: Pre = 0.8481 -> Post = 0.8319, Self-Shift = 0.2621
  Plurality (CovD): Pre = 0.2806 -> Post = 0.3079
  DQI (Total = 14.83):
    - Level of Justification: 3.00
    - Content of Justification: 3.00
    - Respect: 3.00
    - Constructive Politics: 3.00
    - Interactivity: 2.83

Run B (Treatment):
  Convergence: Pre = 0.8481 -> Post = 0.8362, Self-Shift = 0.2225
  Plurality (CovD): Pre = 0.2806 -> Post = 0.3008
  DQI (Total = 14.67):
    - Level of Justification: 3.00
    - Content of Justification: 3.00
    - Respect: 3.00
    - Constructive Politics: 3.00
    - Interactivity: 2.67
\end{verbatim}
\vspace{0.5em}
\textbf{Conclusion:}\\
This was a pilot ($n=1$ run per condition) without a baseline noise floor. The differences between the Control and Treatment are indistinguishable from sampling noise, thus we find no detectable effect of the treatment in this pilot. The null hypothesis ($H_0$) is retained.
\end{tcolorbox}

\subsection{Takeaways}
The addition of the internal reflection step yielded fascinating dynamics in the latest run, but as a pilot with $n=1$, these observations are indistinguishable from sampling noise:
\begin{itemize}
    \item \textbf{More Stubbornness (Lower Self-Shift)}: Treatment agents were actually \textit{more stubborn}. The \verb|<private_scratchpad>| may have served to anchor the agents in their own arguments rather than facilitating empathy (Self-Shift: 0.2225 in Treatment vs 0.2621 in Control). 
    \item \textbf{Plurality Expansion (CovD)}: Interestingly, in \textit{both} runs, the plurality (CovD) increased post-deliberation (Control: 0.2806 $\rightarrow$ 0.3079, Treatment: 0.2806 $\rightarrow$ 0.3008). This indicates that rather than reaching a compromise, the deliberation actually caused the agents' stances to polarize and span a wider volume of opinion space than when they started.
    \item \textbf{Methodological Confounds}: Beyond sampling noise, several confounds affect the treatment run. Firstly, the post-deliberation stance generation does not restate the original persona, risking persona decay over the context window. Secondly, the topic prompt explicitly asks for consensus, while the personas are ``strongly" opposed, meaning convergence partly measures instruction-following against persona conflict rather than genuine opinion change. Lastly, we do not know if the internal reflection changes the response length or style. 
\end{itemize}

Overall, without running multiple seeds with confidence intervals and a noise-floor control (e.g., regenerating the same stance twice at temperature 0.7 to establish baseline Self-Shift), these results remain speculative. We find no detectable effect in this pilot, and the null hypothesis ($H_0$) is retained.

\subsection{Immediate Observations}
\begin{itemize}
    \item Cosine similarity as a metric may not be nuanced enough. As can be observed, the Convergence figures are extremely high and similar ($\sim$0.83) across both runs but this may be a result of the agents sharing a lot of trivially identical phrases (``Thank you...", ``On the topic of...", "I appreciate...")\footnote{This is a shared limitation to \textbf{Park et al. 2023} who remarked that their simulated persona agents were also overly formal and repetitive. They speculate this to be a result of the Instruction Tuning of the underlying LLMs to be suitable chatbots \& practice conversational techniques such as mirroring}.
    \item The above is reinforced by that Plurality of content across stances is very low (0.3079 at best).
    \item Currently, due to the stateless nature of the API calls to the external service, as the deliberation progresses the agents have to be fed with the topic and entire past chat history everytime they are queried. The agents then consider the context when responding. It is unclear whether this has an effect on the deliberation.
    \item LLM-as-a-Judge of DQI is extremely lenient, giving both runs a near-perfect and identical DQI of 14.83/15. This is even with feeding the Judge a complete operational definition of DQI's composite metrics, and examples of how each metric should be graded for this specific topic\footnote{For reference, in the original DQI paper which enlisted 2 human coders to grade a parliamentary session, the average DQI was 7/15.}. This may either be a result of the Judge lacking the ability to differentiate finely between the levels/content of justification between the agents (i.e. the model is not tuned to this task) or the agents' overly formal style of discussion misleading the Judge into rating the discourse quality very highly (i.e. the Judge may have set itself a low bar).
    \item There is an internal tension in the DQI evaluation regarding context. The system prompt explicitly instructs the Judge to score ``based ONLY on the provided turn". While this turn-isolated format is how DQI is scored in the original paper, it neglects that the human judges of DQI still have memory of the earlier utterances in the deliberations which influence their scores.  Especially for relational criteria like Interactivity and Respect. 
\end{itemize}

\newpage{}

\section{Multi-Trial Replication Results (\texorpdfstring{$N = 30$}{N = 30})}
\label{sec:replicationresults}
\rule{\textwidth}{0.4pt}

\subsection{Replication Overview and Dataset}
To establish statistical significance beyond the single-seed pilot ($n=1$), a pre-registered replication suite consisting of $N = 30$ independent paired trials (60 complete deliberation runs, 360 dialogue turns) was executed. Each trial generated a shared pre-deliberation baseline stance at temperature $T = 0.70$, followed by parallel Control (direct dialogue response) and Treatment (mandatory \texttt{private\_scratchpad} internal monologue) deliberation runs.\\

The entire 30-trial replication dataset has been exported, versioned, and hosted on \href{https://huggingface.co/datasets/yashpola/simple-simulation-school-closure}{Hugging Face}. The repository provides two structured subsets in Parquet and JSONL formats:
\begin{itemize}
    \item \textbf{\texttt{trials}}: 60 rows (30 Control + 30 Treatment) recording trial identifiers, condition, semantic convergence, self-shift, pre- and post-deliberation CovD, DQI sub-scores, and baseline/final stances.
    \item \textbf{\texttt{turns}}: 360 rows recording each conversational turn, speaker persona, round number, internal scratchpad reflections, and public utterances.
\end{itemize}

\subsection{Statistical Evaluation}
All indicators are reported as $\text{Mean} \pm \text{Standard Deviation}$ across the 30 independent paired trials. Differences are computed within pairs ($\Delta = \text{Treatment} - \text{Control}$), evaluated using paired two-tailed $t$-tests and Wilcoxon signed-rank tests.

\begin{table}[H]
\centering
\caption{Primary Metrics across $N = 30$ Deliberation Trials}
\label{tab:replication_primary}
\vspace{0.5em}
\resizebox{\textwidth}{!}{%
\begin{tabular}{lcccccc}
\toprule
\textbf{Metric} & \textbf{Pre-Debate} & \textbf{Control (Direct)} & \textbf{Treatment (Scratchpad)} & \textbf{Difference ($\Delta$)} & \textbf{95\% CI} & \textbf{Paired $t$-test} \\
\midrule
\textbf{Convergence} & $0.8591 \pm 0.0302$ & $0.8758 \pm 0.0413$ & $0.8819 \pm 0.0303$ & $+0.0061 \pm 0.0517$ & $[-0.0124, +0.0247]$ & $t = 0.6501, p = 0.5207$ \\
\textbf{Self-Shift} & --- & $0.1739 \pm 0.0415$ & $0.1814 \pm 0.0402$ & $+0.0075 \pm 0.0607$ & $[-0.0142, +0.0292]$ & $t = 0.6760, p = 0.5044$ \\
\textbf{Plurality (CovD)} & $0.2611 \pm 0.0516$ & $0.2313 \pm 0.0713$ & $0.2213 \pm 0.0525$ & $-0.0100 \pm 0.0895$ & $[-0.0421, +0.0220]$ & $t = -0.6139, p = 0.5441$ \\
\textbf{DQI (Total)} & --- & $14.7389 \pm 0.1132$ & $14.7389 \pm 0.0840$ & $0.0000 \pm 0.1384$ & $[-0.0495, +0.0495]$ & $t = 0.0000, p = 1.0000$ \\
\bottomrule
\end{tabular}%
}
\end{table}

\begin{table}[H]
\centering
\caption{DQI Sub-Dimensional Breakdown (Scale: 1.0 to 3.0)}
\label{tab:replication_dqi}
\vspace{0.5em}
\resizebox{0.85\textwidth}{!}{%
\begin{tabular}{lcccc}
\toprule
\textbf{DQI Dimension} & \textbf{Control (Mean $\pm$ SD)} & \textbf{Treatment (Mean $\pm$ SD)} & \textbf{Difference ($\Delta$)} & \textbf{Paired $t$-test} \\
\midrule
\textbf{Level of Justification} & $3.0000 \pm 0.0000$ & $3.0000 \pm 0.0000$ & $0.0000$ & --- \textit{(Ceiling)} \\
\textbf{Content of Justification} & $3.0000 \pm 0.0000$ & $3.0000 \pm 0.0000$ & $0.0000$ & --- \textit{(Ceiling)} \\
\textbf{Respect} & $3.0000 \pm 0.0000$ & $3.0000 \pm 0.0000$ & $0.0000$ & --- \textit{(Ceiling)} \\
\textbf{Constructive Politics} & $2.9889 \pm 0.0609$ & $3.0000 \pm 0.0000$ & $+0.0111$ & $t = 1.0000, p = 0.3256$ \\
\textbf{Interactivity} & $2.7500 \pm 0.0848$ & $2.7389 \pm 0.0840$ & $-0.0111$ & $t = -0.4937, p = 0.6253$ \\
\bottomrule
\end{tabular}%
}
\end{table}

\subsection{Key Findings}
\begin{itemize}
    \item \textbf{Dissipation of Pilot ``Stubbornness'' Effect}: In the single-seed pilot ($n = 1$), the Treatment condition appeared to increase stubbornness (Self-Shift decreased from $0.2621$ to $0.2225$). In the $N = 30$ replication, this discrepancy completely evaporated: Control Self-Shift averaged $0.1739 \pm 0.0415$ and Treatment averaged $0.1814 \pm 0.0402$. The mean difference ($+0.0075$) is statistically indistinguishable from zero ($p = 0.5044$), showing that the pilot result was an artifact of stochastic temperature sampling at $T = 0.70$.
    \item \textbf{Conversational Convergence Without Reflection Effect}: Deliberation consistently induced convergence between agents relative to their pre-deliberation baseline ($0.8591 \rightarrow 0.8758$ in Control, $0.8819$ in Treatment). However, the additional $+0.0061$ gain from internal reflection is statistically insignificant ($p = 0.5207$). Agents moved closer through the interactive conversation itself, not because of the private scratchpad.
    \item \textbf{Contraction of Opinion Space (CovD)}: Contrary to the pilot where CovD expanded post-deliberation ($0.2806 \rightarrow 0.3079$), the 30-trial replication demonstrated systematic opinion domain contraction from baseline ($0.2611 \pm 0.0516$) down to $0.2313 \pm 0.0713$ in Control and $0.2213 \pm 0.0525$ in Treatment. The small difference between conditions ($-0.0100$) was not significant ($p = 0.5441$).
    \item \textbf{DQI Metric Saturation and Ceiling Effects}: The LLM-as-a-Judge assigned identical mean total scores across both conditions ($14.7389 / 15.0000$, $p = 1.0000$). Three dimensions (\textit{Level of Justification}, \textit{Content of Justification}, \textit{Respect}) remained pegged at maximum scores of $3.0000 \pm 0.0000$ across all 360 turns, indicating that \texttt{Llama-3.3-70B} inherently adheres to civility and justification norms, creating a ceiling effect that prevents the metric from detecting fine-grained differences.
    \item \textbf{Scratchpad as Role-Reinforcement Sandbox}: Qualitative analysis of the internal monologues reveals that agents utilized the \texttt{private\_scratchpad} not to reconsider their fundamental stances, but rather to formulate rhetorical packaging for their assigned institutional roles. Consequently, internal reflection acted as a strategic buffer rather than an engine for ideological concession.
\end{itemize}

\subsection{Conclusion}
The empirical results across 30 independent paired trials provide no evidence that mandating an internal reflection step improves deliberative consensus, willingness to compromise, or discourse quality in dyadic citizen panels. The Null Hypothesis ($\boldsymbol{H}_0$) is definitively retained across all evaluated dimensions.

\newpage{}

\section{Codebase Reproducability}
\label{sec:taskdescription}
\rule{\textwidth}{0.4pt}

\subsection{Coding Agent Information}
The repo also contains an updated \texttt{experiment\_design\_school\_closures.md} and \texttt{coding\_guidelines.md} which can help a future agent reproduce the codebase in its current structure.\\

The information in both manifests was manually specified, but written automatically by \texttt{Gemini 3.1 Pro (Thinking-High)}.

\subsection{Tests}
There are 2 test scripts in \texttt{repo/tests}.
\begin{itemize}
    \item The first script verifies that an API key and Base URL is present and correctly formatted. It then pings every LLM model specified in \texttt{config.json} to ensure it is responsive and accessible on the service.
    \item The second script validates the regex parsing logic in \texttt{extract\_tags} to ensure the main script correctly extracts the \texttt{private\_scratchpad} from the model's response.
\end{itemize}

Both tests can be run using: \texttt{python -m unittest discover -s tests}.\\

A sample test result looks like:
\begin{sidequote}
    \begin{verbatim}
    test_api_key_set (test_api_check.TestAPICheck.test_api_key_set) ... ok
    test_models_accessible (test_api_check.TestAPICheck.test_models_accessible) ... ok
    test_extract_tags_both (test_sanity.TestSanityCheck.test_extract_tags_both) ... ok
    test_extract_tags_missing_public (test_sanity.TestSanityCheck.test_extract_tags_missing_public) ... ok
    test_extract_tags_no_tags (test_sanity.TestSanityCheck.test_extract_tags_no_tags) ... ok
    
    ----------------------------------------------------------------------
    Ran 5 tests in 1.986s
    
    OK
    \end{verbatim}
\end{sidequote}

\newpage

\section{References}
\label{sec:taskdescription}
\rule{\textwidth}{0.4pt}\\

Argyle, Lisa P., et al. "Out of one, many: Using language models to simulate human samples." Political Analysis 31.3 (2023): 337-351.\\

Bächtiger, André, et al. "Deliberative democracy and its measures." Oxford Handbook of Deliberative Democracy. Oxford University Press, 2022.\\

Hayati, Shirley Anugrah, et al. "How far can we extract diverse perspectives from large language models?." Proceedings of the 2024 Conference on Empirical Methods in Natural Language Processing. 2024.\\

Ke, Zhisheng, et al. "Applying item response theory and classical test theory to evaluate large language models on standardized assessments." Proceedings of the International Conference on Learning Representations (ICLR). 2026.\\

Nguyen, Vivian, Lillian Lee, and Cristian Danescu-Niculescu-Mizil. "Hanging in the balance: pivotal moments in crisis counseling conversations." Proceedings of the 63rd Annual Meeting of the Association for Computational Linguistics (Volume 1: Long Papers). 2025.\\

Park, Joon Sung, et al. "Generative agents: Interactive simulacra of human behavior." Proceedings of the 36th annual acm symposium on user interface software and technology. 2023.\\

Reimers, Nils, and Iryna Gurevych. "Sentence-bert: Sentence embeddings using siamese bert-networks." Proceedings of the 2019 conference on empirical methods in natural language processing and the 9th international joint conference on natural language processing (EMNLP-IJCNLP). 2019.\\

Steenbergen, Marco R., et al. "Measuring political deliberation: A discourse quality index." Comparative European Politics 1.1 (2003): 21-48.\\

Veri, Francesco. "Replicability in AI-driven deliberation: Measuring latent structural stability over textual identity." Journal of Artificial Societies and Social Simulation (2026).\\

Wang, Angelina, Jamie Morgenstern, and John P. Dickerson. "Large language models that replace human participants can harmfully misportray and flatten identity groups." Nature Machine Intelligence 7.3 (2025): 400-411.\\

Yang, Yi-Heng, and Götz Bachmann. "Deliberation and voting in democratic decision-making: Complementary mechanisms for collective choice." arXiv preprint arXiv:2502.05017 (2025).\\

Zhuang, Kaiyao, Yuheng Wang, and Weiyu Zhang. "Can AI Deliberate? Evaluating Deliberative Quality and Stance Flow in Multi-Agent LLMs." Proceedings of the Extended Abstracts of the 2026 CHI Conference on Human Factors in Computing Systems. 2026.

\end{document}